\bibitem{NCIwaniecConv06} H. Iwaniec. Conversations on the exceptional character. In \emph{Analytic Number Theory} (Cetraro, 2002), Lecture Notes in Mathematics 1891, Springer, 2006, 97--132. % verified: DOI 10.1007/978-3-540-36364-4_3, as supplied by the authors; the text was not read, and we do not know whether it describes the elliptic curve family of Section~1
\bibitem{NCConreyIwaniec02} J. B. Conrey and H. Iwaniec. Spacing of zeros of Hecke $L$-functions and the class number problem. \emph{Acta Arithmetica} 103 (2002), 259--312. % verified: DOI 10.4064/aa103-3-5, as supplied by the authors
\bibitem{NCConrey03} J. B. Conrey. The Riemann hypothesis. \emph{Notices of the American Mathematical Society} 50 (2003), 341--353. % UNCHECKED: fields from Conrey's publication list only; the survey text was not read
\bibitem{NCMestre86} J.-F. Mestre. Formules explicites et minorations de conducteurs de vari\'et\'es alg\'ebriques. \emph{Compositio Mathematica} 58 (1986), 209--232. % verified: Numdam CM_1986__58_2_209_0 and EuDML 89769, 2026-11-04 (no DOI)
\bibitem{NCBrumer92} A. Brumer. The average rank of elliptic curves I. \emph{Inventiones Mathematicae} 109 (1992), 445--472. % verified: CrossRef 10.1007/bf01232033 2026-11-04
\bibitem{NCBoas54} R. P. Boas. \emph{Entire Functions}. Academic Press, New York, 1954. % verified: publisher and year, WorldCat 847696, 2026-11-04; pages 124 to 132 for Akhiezer's factorisation per secondary sources, book not read
\bibitem{NCGrossZagier86} B. H. Gross and D. B. Zagier. Heegner points and derivatives of $L$-series. \emph{Inventiones Mathematicae} 84 (1986), 225--320. % verified: CrossRef 10.1007/BF01388809 2026-11-04
\bibitem{NCKolyvagin89} V. A. Kolyvagin. Finiteness of $E(\mathbb Q)$ and $\mathrm{Sha}(E,\mathbb Q)$ for a subclass of Weil curves. \emph{Mathematics of the USSR Izvestiya} 32 (1989), 523--541. % verified: CrossRef 10.1070/im1989v032n03abeh000779 2026-11-04 (English translation of the 1988 Izvestiya paper)
\bibitem{BohrCourant} H. Bohr and R. Courant. Neue Anwendungen der Theorie der Diophantischen Approximationen auf die Riemannsche Zetafunktion. \emph{Journal f\"ur die reine und angewandte Mathematik} 144 (1914), 249--274. % UNCHECKED: fields from the bibliography of the note, not seen in the original
